\scanpage{44}
% archive foliation: 43

$\mathcal T(\overline{\mathbb F}_p)$ est un torseur canonique sous
\[
\begin{gathered}
H^1(\mathbb A,\mu_2)\\
\parallel\\
\prod_{\ell}\mathbb Z_\ell^*/\mathbb Z_\ell^{*2}\\
\text{(y compris $p$ et $\infty$)}
\end{gathered}
\]

Ce torseur à priori dépend encore de $\overline{\mathbb F}_p$, \deleted{\unsure{on doit donc}} et cette dépendance est précisée par la connaissance de l'opération de $\widehat{\mathbb Z}\ni f$ sur $\mathcal T(\overline{\mathbb F}_p)$.
Cette opération se fait via les opérations de $\widehat{\mathbb Z}$ sur $C(A',T')$,
\inserted{(opère \unsure{habituellement} sur $A'(1)$, $f$ opère par $p$) et}
\deleted{\unsure{$C(A_p,T_p)$}}
\inserted{$C(A_p,T_p)$}
(via opération de $\deleted{\unsure{\mathbb Z}}\mathbb Zf$ sur $C(A_p,T_p)$ par Frobenius) et l'opération triviale sur $C(A_\infty,T_\infty)$.
Donc on trouve que $f\in\widehat{\mathbb Z}$ opère par le même élément $(p,f,p)$ que plus haut, donc
\[
f\text{ opère trivialement sur }\mathcal T(\overline{\mathbb F}_p)
\]
Ainsi, on a un torseur canonique \inserted{$\mathcal T$} sous $H^1(\mathbb A,\mu_2)$.

\uline{NB} Ceci suggère qu'on a une \deleted{\unsure{gerbe de catégories}} \inserted{gerbe triviale (ou} \unsure{trivialisée)} de \deleted{\unsure{faisceaux}} $\mu_2$ sur $\mathbb A$, dont $\mathcal T$ serait le torseur des classes d'isomorphie.
On peut en effet considérer les couples $(C,\deleted{\unsure{\xi}}\alpha)$, où $C$ est une catégorie de Deuring de type $\xi\in\mathrm{Br}(\mathbb Q)$, \inserted{d'où $T_C=T$,} et $\alpha$ une \deleted{\unsure{isom}} classe d'équivalences de classes d'isomorphie \inserted{\unsure{le privilège}} de $(C,T)_{\mathbb Z_p}$ avec $C(\mathbb Z_p),T(\mathbb Z_p)$ modulo équivalences provenant de $\uline{\operatorname{Pic}}(\mathbb Z_p)$; i.e. une section du torseur de groupe $\mathbb Z/2\mathbb Z$ déduit de celui de groupe
\[
\unsure{\mathrm{Eq}}(C(\mathbb Z_p),T(\mathbb Z_p))
=\mathbb Z_p^*/\mathbb Z_p^{*2}\times\mathbb Z/2\mathbb Z
\]
\unsure{relatif} par \unsure{extension} au groupe \unsure{standard} à l'aide de $\operatorname{pr}_2$.

\margin{\deleted{\unsure{à isom près}}
On \unsure{précise}: (objets: \unsure{quotients} $(M,\sigma)$): \uline{isom} \uline{près}
(il y a deux classes d'isomorphie possibles; si $(M,\sigma)$ est l'une d'elles, l'autre est $(M^{(p)},\sigma)$ \unsure{qui} \unsure{s'en déduit en appliquant} Frobenius).}

\margin{\unsure{plus compliqué}, il suffit de \unsure{fixer une} classe d'isom. près d'équivalences de cat.
$C_{\mathbb Z_p}\simeq C(\mathbb Z_p)$ \unsure{de Deuring}.}
\margin{en termes concrets.
un $\alpha$ est donné ($C$ défini par $O$) en se donnant une \unsure{isom} de $O\otimes_{\mathbb Z}\mathbb Z_p$ avec l'anneau des endom. de $\unsure{M_0}$; où \unsure{$M_0$ est un} cristal sur $\overline{\mathbb F}_p$ (\unsure{de rang $2$}),
\deleted{\unsure{on dispose}\ \unsure{$\widehat A/\widehat A^{\,2}$}\ \illegible{3}}.}

Alors tout $C$ admet un $\alpha$, et deux $(C,\alpha)$ $(C',\alpha')$
